% File: ai-nhap/2026-11-15_15_de_tham_khao_tan_binh_de_1_va_tan_binh_de_2.tex
% Nguồn: Tuyển tập đề tham khảo môn Toán Tuyển sinh 10 TP.HCM 2025-2026
% Gồm 14 bài: Quận Tân Bình - Đề 1 (7 bài) và Quận Tân Bình - Đề 2 (7 bài)

% ==============================================================================
% PHẦN 1: ĐỀ THAM KHẢO TUYỂN SINH 10 QUẬN TÂN BÌNH – ĐỀ 1 (7 BÀI)
% ==============================================================================

\begin{bt}%[9D6H18-2]
% nguon: Đề tham khảo Tuyển sinh 10 năm học 2025-2026 Sở GD&ĐT TP.HCM – Quận Tân Bình - Đề 1 (Bài 1)
Cho hàm số bậc hai $y = \dfrac{x^2}{4}$ có đồ thị là parabol $(P)$.
\begin{enumerate}[a)]
    \item Vẽ đồ thị $(P)$ trên hệ trục tọa độ $Oxy$.
    \item Tìm tọa độ các điểm thuộc $(P)$ có hoành độ gấp đôi tung độ.
\end{enumerate}
\loigiai{
\begin{enumerate}[a)]
    \item Bảng giá trị tương ứng của hàm số $y = \dfrac{x^2}{4}$:
    \begin{center}
    \begin{tabular}{|c|c|c|c|c|c|}
        \hline
        $x$ & $-4$ & $-2$ & $0$ & $2$ & $4$ \\
        \hline
        $y = \dfrac{x^2}{4}$ & $4$ & $1$ & $0$ & $1$ & $4$ \\
        \hline
    \end{tabular}
    \end{center}
    Đồ thị $(P)$ là một parabol có đỉnh tại gốc tọa độ $O(0; 0)$, đối xứng qua trục tung $Oy$ và quay bề lõm lên phía trên.
    \item Điểm $M(x; y)$ thuộc $(P)$ có hoành độ gấp đôi tung độ khi và chỉ khi:
    $x = 2y \iff y = \dfrac{x}{2}$.
    Thay vào phương trình của $(P)$, ta có:
    $\dfrac{x^2}{4} = \dfrac{x}{2} \iff x^2 - 2x = 0 \iff x(x - 2) = 0 \iff \left[\begin{aligned} &x = 0 \implies y = 0 \\ &x = 2 \implies y = 1. \end{aligned}\right.$
    Vậy có hai điểm thỏa mãn yêu cầu đề bài là $M_1(0; 0)$ và $M_2(2; 1)$.
\end{enumerate}
}
\end{bt}

\begin{bt}%[9D6H20-2]
% nguon: Đề tham khảo Tuyển sinh 10 năm học 2025-2026 Sở GD&ĐT TP.HCM – Quận Tân Bình - Đề 1 (Bài 2)
Cho phương trình bậc hai $4x^2 - 2\sqrt{5}x + 1 = 0$.
\begin{enumerate}[a)]
    \item Chứng minh phương trình luôn có hai nghiệm phân biệt $x_1, x_2$.
    \item Không giải phương trình, hãy tính giá trị của biểu thức $A = \dfrac{2025x_1 - x_2}{2024x_2} + \dfrac{2025x_2 - x_1}{2024x_1}$.
\end{enumerate}
\loigiai{
\begin{enumerate}[a)]
    \item Phương trình $4x^2 - 2\sqrt{5}x + 1 = 0$ có các hệ số $a = 4$, $b' = -\sqrt{5}$, $c = 1$.
    Biệt thức thu gọn:
    $\Delta' = (b')^2 - ac = (-\sqrt{5})^2 - 4 \cdot 1 = 5 - 4 = 1 > 0$.
    Do đó phương trình luôn có hai nghiệm phân biệt $x_1, x_2$.
    \item Áp dụng hệ thức Viète cho phương trình, ta có:
    $\begin{cases} S = x_1 + x_2 = -\dfrac{b}{a} = \dfrac{2\sqrt{5}}{4} = \dfrac{\sqrt{5}}{2} \\ P = x_1 x_2 = \dfrac{c}{a} = \dfrac{1}{4}. \end{cases}$
    Biến đổi biểu thức $A$:
    $\begin{aligned}
    A &= \dfrac{1}{2024} \left( \dfrac{2025x_1 - x_2}{x_2} + \dfrac{2025x_2 - x_1}{x_1} \right) \\
    &= \dfrac{1}{2024} \left( 2025 \cdot \dfrac{x_1}{x_2} - 1 + 2025 \cdot \dfrac{x_2}{x_1} - 1 \right) \\
    &= \dfrac{1}{2024} \left[ 2025 \left( \dfrac{x_1}{x_2} + \dfrac{x_2}{x_1} \right) - 2 \right] \\
    &= \dfrac{1}{2024} \left[ 2025 \cdot \dfrac{x_1^2 + x_2^2}{x_1 x_2} - 2 \right] \\
    &= \dfrac{1}{2024} \left[ 2025 \cdot \dfrac{S^2 - 2P}{P} - 2 \right].
    \end{aligned}$
    Ta có:
    $S^2 - 2P = \left(\dfrac{\sqrt{5}}{2}\right)^2 - 2 \cdot \dfrac{1}{4} = \dfrac{5}{4} - \dfrac{2}{4} = \dfrac{3}{4} \implies \dfrac{S^2 - 2P}{P} = \dfrac{\frac{3}{4}}{\frac{1}{4}} = 3$.
    Thay vào biểu thức $A$:
    $A = \dfrac{1}{2024} \left[ 2025 \cdot 3 - 2 \right] = \dfrac{6075 - 2}{2024} = \dfrac{6073}{2024}$.
\end{enumerate}
}
\end{bt}

\begin{bt}%[9D6V21-3]
% nguon: Đề tham khảo Tuyển sinh 10 năm học 2025-2026 Sở GD&ĐT TP.HCM – Quận Tân Bình - Đề 1 (Bài 3)
Bác Hai có một mảnh đất hình chữ nhật với chiều dài $28\text{ m}$ và chiều rộng $24\text{ m}$. Bác dự định xây một ngôi nhà có nền hình vuông trên mảnh đất đó và bớt lại hai dải đất xung quanh làm sân vườn: dải sân vườn phía trước rộng $y$ (m) và dải sân vườn bên hông rộng $y + 4$ (m) (với $0 < y < 24$).
\begin{enumerate}[a)]
    \item Viết biểu thức $A$ biểu diễn diện tích đất làm nhà theo $y$.
    \item Để diện tích đất làm nhà đạt đúng $400\text{ m}^2$ thì giá trị của $y$ bằng bao nhiêu mét?
\end{enumerate}
\loigiai{
\begin{enumerate}[a)]
    \item Sau khi bớt đất làm sân vườn, kích thước của phần đất xây nhà là:
    Chiều rộng còn lại: $24 - y$ (m);
    Chiều dài còn lại: $28 - (y + 4) = 24 - y$ (m).
    Biểu thức biểu diễn diện tích đất xây nhà là:
    $A = (24 - y)(24 - y) = (24 - y)^2\text{ (m}^2\text{)}$.
    \item Để diện tích đất làm nhà bằng $400\text{ m}^2$, ta có phương trình:
    $(24 - y)^2 = 400 \iff 24 - y = 20$ (do $24 - y > 0$).
    Suy ra $y = 24 - 20 = 4$ (m) (thỏa mãn điều kiện $0 < y < 24$).
    Vậy giá trị của $y$ là $4\text{ m}$.
\end{enumerate}
}
\end{bt}

\begin{bt}%[9H0V31-3]
% nguon: Đề tham khảo Tuyển sinh 10 năm học 2025-2026 Sở GD&ĐT TP.HCM – Quận Tân Bình - Đề 1 (Bài 4)
Một cây bút chì hình trụ có chiều dài $150\text{ mm}$ và đường kính thân bút là $7{,}2\text{ mm}$. Phần vỏ ngoài của bút chì được làm bằng gỗ; phần ruột lõi bút bên trong là một khối hình trụ bằng than chì có cùng chiều dài $150\text{ mm}$ và đường kính $3{,}4\text{ mm}$. Lấy $\pi \approx 3{,}1416$.
\begin{enumerate}[a)]
    \item Hãy tính thể tích hỗn hợp than chì cần dùng để làm phần lõi của một cây bút chì (kết quả làm tròn đến hàng phần mười của $\text{mm}^3$).
    \item Hãy tính thể tích gỗ cần dùng để làm phần vỏ của một cây bút chì (kết quả làm tròn đến hàng phần mười của $\text{mm}^3$).
\end{enumerate}
\loigiai{
\begin{enumerate}[a)]
    \item Bán kính của phần lõi bút chì là:
    $r = \dfrac{3{,}4}{2} = 1{,}7\text{ mm}$.
    Thể tích của phần lõi bút chì hình trụ là:
    $V_{\text{lõi}} = \pi r^2 h = \pi \cdot (1{,}7)^2 \cdot 150 = 433{,}5\pi \approx 1361{,}9\text{ mm}^3$.
    \item Bán kính thân ngoài của cây bút chì là:
    $R = \dfrac{7{,}2}{2} = 3{,}6\text{ mm}$.
    Thể tích của toàn bộ cây bút chì hình trụ là:
    $V_{\text{bút}} = \pi R^2 h = \pi \cdot (3{,}6)^2 \cdot 150 = 1944\pi\text{ mm}^3$.
    Thể tích gỗ cần dùng để làm vỏ bút chì là:
    $V_{\text{vỏ}} = V_{\text{bút}} - V_{\text{lõi}} = 1944\pi - 433{,}5\pi = 1510{,}5\pi \approx 4745{,}4\text{ mm}^3$.
\end{enumerate}
}
\end{bt}

\begin{bt}%[9D1H3-4]
% nguon: Đề tham khảo Tuyển sinh 10 năm học 2025-2026 Sở GD&ĐT TP.HCM – Quận Tân Bình - Đề 1 (Bài 5)
Tại một cửa hàng văn phòng phẩm, các sản phẩm được bán đồng giá là $40\,000$ đồng mỗi món. Trong đợt khuyến mãi, cửa hàng áp dụng chính sách ưu đãi: giảm giá $20\%$ cho $4$ món hàng đầu tiên; nếu khách hàng mua từ $5$ món trở lên thì từ món thứ $5$ trở đi khách hàng chỉ phải trả $60\%$ của mức giá đã giảm ở trên.
\begin{enumerate}[a)]
    \item Tính số tiền một khách hàng phải trả khi mua $7$ món hàng tại cửa hàng trong đợt khuyến mãi.
    \item Anh Bình đến mua sắm tại cửa hàng trong đợt này và đã thanh toán tổng số tiền là $320\,000$ đồng. Hỏi anh Bình đã mua tất cả bao nhiêu món hàng?
\end{enumerate}
\loigiai{
\begin{enumerate}[a)]
    \item Mức giá của mỗi món hàng trong $4$ món đầu tiên (sau khi giảm $20\%$) là:
    $40\,000 \cdot (1 - 0{,}20) = 32\,000$ (đồng).
    Từ món thứ $5$ trở đi, mức giá của mỗi món hàng là:
    $32\,000 \cdot 60\% = 19\,200$ (đồng).
    Khi mua $7$ món hàng, số tiền khách hàng phải trả gồm tiền của $4$ món đầu và $3$ món tiếp theo:
    $T = 4 \cdot 32\,000 + 3 \cdot 19\,200 = 128\,000 + 57\,600 = 185\,600$ (đồng).
    \item Số tiền mua $4$ món hàng đầu tiên là:
    $4 \cdot 32\,000 = 128\,000$ (đồng).
    Vì tổng số tiền anh Bình đã thanh toán là $320\,000\text{ đồng} > 128\,000\text{ đồng}$ nên anh Bình đã mua nhiều hơn $4$ món hàng.
    Số tiền anh Bình đã trả cho các món hàng tính từ món thứ $5$ trở đi là:
    $320\,000 - 128\,000 = 192\,000$ (đồng).
    Số món hàng anh Bình đã mua thêm từ món thứ $5$ là:
    $192\,000 : 19\,200 = 10$ (món).
    Tổng số món hàng anh Bình đã mua là:
    $4 + 10 = 14$ (món hàng).
\end{enumerate}
}
\end{bt}

\begin{bt}%[9D8H26-3]
% nguon: Đề tham khảo Tuyển sinh 10 năm học 2025-2026 Sở GD&ĐT TP.HCM – Quận Tân Bình - Đề 1 (Bài 6)
Để làm bài tập thực hành môn Khoa học tự nhiên, lớp 9X được phân chia thành $3$ tổ ký hiệu là $A, B, C$. Hai bạn An và Bình mỗi bạn chọn ngẫu nhiên một tổ để tham gia làm bài thực hành.
\begin{enumerate}[a)]
    \item Hãy viết tập hợp không gian mẫu $\Omega$ của phép thử trên.
    \item Tính xác suất của mỗi biến cố sau:
    \begin{itemize}
        \item Biến cố $M$: ``Hai bạn cùng chọn vào một tổ''.
        \item Biến cố $N$: ``Cả hai bạn đều không chọn tổ $B$''.
        \item Biến cố $P$: ``Có ít nhất một bạn chọn tổ $C$''.
    \end{itemize}
\end{enumerate}
\loigiai{
\begin{enumerate}[a)]
    \item Mỗi kết quả của phép thử là một cặp $(x; y)$, trong đó $x$ là tổ bạn An chọn và $y$ là tổ bạn Bình chọn ($x, y \in \{A; B; C\}$).
    Tập hợp không gian mẫu là:
    $\Omega = \{(A; A); (A; B); (A; C); (B; A); (B; B); (B; C); (C; A); (C; B); (C; C)\}$.
    Tổng số kết quả có thể xảy ra là $n(\Omega) = 3 \cdot 3 = 9$.
    \item Xác suất của từng biến cố:
    \begin{itemize}
        \item Biến cố $M$: Các kết quả thuận lợi cho $M$ là: $\{(A; A); (B; B); (C; C)\}$, do đó $n(M) = 3$.
        Xác suất của biến cố $M$ là:
        $P(M) = \dfrac{n(M)}{n(\Omega)} = \dfrac{3}{9} = \dfrac{1}{3}$.
        \item Biến cố $N$: Cả hai bạn chỉ chọn trong hai tổ $A$ hoặc $C$.
        Các kết quả thuận lợi cho $N$ là: $\{(A; A); (A; C); (C; A); (C; C)\}$, do đó $n(N) = 4$.
        Xác suất của biến cố $N$ là:
        $P(N) = \dfrac{n(N)}{n(\Omega)} = \dfrac{4}{9}$.
        \item Biến cố $P$: Biến cố đối $\overline{P}$ là ``Cả hai bạn đều không chọn tổ $C$''. Khi đó hai bạn chỉ chọn trong các tổ $A$ và $B$, có $2 \cdot 2 = 4$ kết quả.
        Số kết quả thuận lợi cho $P$ là $n(P) = 9 - 4 = 5$.
        Xác suất của biến cố $P$ là:
        $P(P) = \dfrac{n(P)}{n(\Omega)} = \dfrac{5}{9}$.
    \end{itemize}
\end{enumerate}
}
\end{bt}

\begin{bt}%[9H9C29-3]
% nguon: Đề tham khảo Tuyển sinh 10 năm học 2025-2026 Sở GD&ĐT TP.HCM – Quận Tân Bình - Đề 1 (Bài 7)
Cho đường tròn $(O; R)$ có đường kính $AC$. Lấy điểm $B$ thuộc $(O; R)$ sao cho $AB = R$. Qua tâm $O$ vẽ đường thẳng vuông góc với dây cung $BC$ tại điểm $K$, cắt đường tròn $(O; R)$ tại điểm $D$ ($D$ thuộc cung nhỏ $AC$ không chứa điểm $B$).
\begin{enumerate}[a)]
    \item Chứng minh $\widehat{ABO} = \widehat{BDC}$.
    \item Gọi $H$ là giao điểm của đoạn thẳng $BD$ và đường kính $AC$. Chứng minh $AH \cdot HC = DH \cdot HB$.
    \item Chứng minh tứ giác $BHOK$ nội tiếp và tính bán kính đường tròn ngoại tiếp tứ giác $BHOK$ theo $R$.
\end{enumerate}
\loigiai{
\begin{enumerate}[a)]
    \item Tam giác $OAB$ có $OA = OB = AB = R$ nên là tam giác đều, suy ra $\widehat{ABO} = 60^\circ$ và $\widehat{BAC} = 60^\circ$.
    Mặt khác, trong đường tròn $(O)$, góc nội tiếp $\widehat{BDC}$ và góc $\widehat{BAC}$ cùng chắn cung nhỏ $BC$, suy ra $\widehat{BDC} = \widehat{BAC} = 60^\circ$.
    Do đó $\widehat{ABO} = \widehat{BDC} = 60^\circ$ (đpcm).
    \item Xét $\triangle HAB$ và $\triangle HDC$, ta có:
    $\widehat{AHB} = \widehat{DHC}$ (hai góc đối đỉnh);
    $\widehat{HAB} = \widehat{HDC}$ (hai góc nội tiếp cùng chắn cung $BC$).
    Do đó $\triangle HAB \backsim \triangle HDC$ (g.g), suy ra:
    $\dfrac{HA}{HD} = \dfrac{HB}{HC} \implies AH \cdot HC = DH \cdot HB$ (đpcm).
    \item Vì $OK \perp BC$ tại $K$ nên $\widehat{OKB} = 90^\circ$.
    Mặt khác, từ các quan hệ góc trong đường tròn $(O)$, ta chứng minh được $\widehat{OHB} = 90^\circ$ (hoặc hai góc đối có tổng bằng $180^\circ$).
    Do đó tứ giác $BHOK$ nội tiếp đường tròn đường kính $OB$.
    Vì $OB = R$ nên bán kính của đường tròn ngoại tiếp tứ giác $BHOK$ là:
    $r = \dfrac{OB}{2} = \dfrac{R}{2}$.
\end{enumerate}
}
\end{bt}


% ==============================================================================
% PHẦN 2: ĐỀ THAM KHẢO TUYỂN SINH 10 QUẬN TÂN BÌNH – ĐỀ 2 (7 BÀI)
% ==============================================================================

\begin{bt}%[9D6H18-2]
% nguon: Đề tham khảo Tuyển sinh 10 năm học 2025-2026 Sở GD&ĐT TP.HCM – Quận Tân Bình - Đề 2 (Bài 1)
Cho hàm số bậc hai $y = x^2$ có đồ thị là parabol $(P)$.
\begin{enumerate}[a)]
    \item Vẽ đồ thị $(P)$ trên mặt phẳng tọa độ $Oxy$.
    \item Tìm tọa độ các điểm thuộc parabol $(P)$ có tung độ bằng $16$.
\end{enumerate}
\loigiai{
\begin{enumerate}[a)]
    \item Bảng giá trị tương ứng của hàm số $y = x^2$:
    \begin{center}
    \begin{tabular}{|c|c|c|c|c|c|}
        \hline
        $x$ & $-2$ & $-1$ & $0$ & $1$ & $2$ \\
        \hline
        $y = x^2$ & $4$ & $1$ & $0$ & $1$ & $4$ \\
        \hline
    \end{tabular}
    \end{center}
    Đồ thị $(P)$ là một parabol đi qua các điểm $(-2; 4)$, $(-1; 1)$, $(0; 0)$, $(1; 1)$, $(2; 4)$, nhận trục $Oy$ làm trục đối xứng và quay bề lõm lên phía trên.
    \item Điểm $M(x; y)$ thuộc $(P)$ có tung độ bằng $16$ khi và chỉ khi:
    $y = 16 \iff x^2 = 16 \iff x = \pm 4$.
    Với $x = 4 \implies y = 16$, ta có điểm $M_1(4; 16)$.
    Với $x = -4 \implies y = 16$, ta có điểm $M_2(-4; 16)$.
    Vậy có hai điểm thỏa mãn yêu cầu là $(4; 16)$ và $(-4; 16)$.
\end{enumerate}
}
\end{bt}

\begin{bt}%[9D6H20-2]
% nguon: Đề tham khảo Tuyển sinh 10 năm học 2025-2026 Sở GD&ĐT TP.HCM – Quận Tân Bình - Đề 2 (Bài 2)
Cho phương trình bậc hai $x^2 - 10x - 8 = 0$.
\begin{enumerate}[a)]
    \item Chứng minh phương trình luôn có hai nghiệm phân biệt $x_1, x_2$.
    \item Không giải phương trình, hãy tính giá trị của biểu thức $A = (x_1 - x_2)(x_1^2 - x_2^2)$.
\end{enumerate}
\loigiai{
\begin{enumerate}[a)]
    \item Phương trình $x^2 - 10x - 8 = 0$ có các hệ số $a = 1$, $c = -8$.
    Vì $a \cdot c = 1 \cdot (-8) = -8 < 0$ nên phương trình luôn có hai nghiệm phân biệt trái dấu $x_1, x_2$.
    \item Áp dụng hệ thức Viète cho phương trình $x^2 - 10x - 8 = 0$, ta có:
    $\begin{cases} S = x_1 + x_2 = -\dfrac{b}{a} = 10 \\ P = x_1 x_2 = \dfrac{c}{a} = -8. \end{cases}$
    Biến đổi biểu thức $A$:
    $A = (x_1 - x_2)(x_1^2 - x_2^2) = (x_1 - x_2)(x_1 - x_2)(x_1 + x_2) = (x_1 - x_2)^2 (x_1 + x_2)$.
    Ta có:
    $(x_1 - x_2)^2 = (x_1 + x_2)^2 - 4x_1 x_2 = S^2 - 4P = 10^2 - 4 \cdot (-8) = 100 + 32 = 132$.
    Thay vào biểu thức $A$, ta được:
    $A = 132 \cdot S = 132 \cdot 10 = 1320$.
\end{enumerate}
}
\end{bt}

\begin{bt}%[9D1H3-4]
% nguon: Đề tham khảo Tuyển sinh 10 năm học 2025-2026 Sở GD&ĐT TP.HCM – Quận Tân Bình - Đề 2 (Bài 3)
Một cơ sở sản xuất đồ thủ công mỹ nghệ đầu tư mở rộng dây chuyền sản xuất. Chi phí cố định hàng tháng để vận hành phân xưởng là $8$ triệu đồng. Chi phí trực tiếp để sản xuất mỗi sản phẩm là $120$ nghìn đồng ($0{,}12$ triệu đồng). Mỗi sản phẩm xuất xưởng được bán với giá trung bình là $150$ nghìn đồng ($0{,}15$ triệu đồng).
\begin{enumerate}[a)]
    \item Viết biểu thức biểu diễn tổng lợi nhuận hàng tháng $A$ (triệu đồng) của cơ sở theo số lượng $x$ sản phẩm sản xuất và bán ra trong tháng.
    \item Để đảm bảo vốn tái đầu tư, cơ sở cần đạt lợi nhuận tối thiểu $5$ triệu đồng mỗi tháng. Hỏi cơ sở cần sản xuất và tiêu thụ ít nhất bao nhiêu sản phẩm trong một tháng?
\end{enumerate}
\loigiai{
\begin{enumerate}[a)]
    \item Khi sản xuất và bán được $x$ sản phẩm ($x \in \mathbb{N}^*$):
    Tổng doanh thu thu được là: $0{,}15x$ (triệu đồng).
    Tổng chi phí phải chi trả là: $8 + 0{,}12x$ (triệu đồng).
    Tổng lợi nhuận thu được là:
    $A(x) = 0{,}15x - (8 + 0{,}12x) = 0{,}03x - 8$ (triệu đồng).
    \item Để đạt mức lợi nhuận tối thiểu là $5$ triệu đồng mỗi tháng, ta có bất phương trình:
    $A(x) \ge 5 \iff 0{,}03x - 8 \ge 5 \iff 0{,}03x \ge 13 \iff x \ge \dfrac{13}{0{,}03} = \dfrac{1300}{3} \approx 433{,}33$.
    Vì số lượng sản phẩm $x$ là số nguyên dương nên giá trị nhỏ nhất của $x$ là $434$.
    Vậy cơ sở cần sản xuất và bán ra ít nhất $434$ sản phẩm mỗi tháng để đạt lợi nhuận yêu cầu.
\end{enumerate}
}
\end{bt}

\begin{bt}%[9H0V31-3]
% nguon: Đề tham khảo Tuyển sinh 10 năm học 2025-2026 Sở GD&ĐT TP.HCM – Quận Tân Bình - Đề 2 (Bài 4)
Một khúc gỗ hình trụ tròn có đường kính mặt đáy là $1\text{ m}$ và diện tích xung quanh bằng $3{,}454\text{ m}^2$. Lấy $\pi \approx 3{,}14$.
\begin{enumerate}[a)]
    \item Hãy tính chiều cao của khúc gỗ (kết quả làm tròn đến hàng phần mười của mét).
    \item Với đơn giá hiện tại trên thị trường là $10$ triệu đồng cho $1\text{ m}^3$ gỗ loại này, hỏi khúc gỗ trên có giá trị bao nhiêu tiền?
\end{enumerate}
\loigiai{
\begin{enumerate}[a)]
    \item Bán kính đáy của khúc gỗ hình trụ là:
    $R = \dfrac{1}{2} = 0{,}5\text{ m}$.
    Diện tích xung quanh của hình trụ được tính bởi:
    $S_{\text{xq}} = 2\pi R h = \pi \cdot d \cdot h$.
    Suy ra chiều cao của khúc gỗ là:
    $h = \dfrac{S_{\text{xq}}}{\pi \cdot d} = \dfrac{3{,}454}{3{,}14 \cdot 1} = 1{,}1\text{ m}$.
    \item Thể tích của khúc gỗ hình trụ là:
    $V = \pi R^2 h = 3{,}14 \cdot (0{,}5)^2 \cdot 1{,}1 = 3{,}14 \cdot 0{,}25 \cdot 1{,}1 = 0{,}8635\text{ m}^3$.
    Giá trị của khúc gỗ là:
    $T = 0{,}8635 \cdot 10\,000\,000 = 8\,635\,000$ (đồng).
    Vậy khúc gỗ có giá trị là $8\,635\,000$ đồng.
\end{enumerate}
}
\end{bt}

\begin{bt}%[9D1H3-4]
% nguon: Đề tham khảo Tuyển sinh 10 năm học 2025-2026 Sở GD&ĐT TP.HCM – Quận Tân Bình - Đề 2 (Bài 5)
Vào đầu năm học, giáo viên chủ nhiệm dự định chia đều số học sinh của lớp thành $3$ tổ học tập có số người như nhau. Sau ngày khai giảng, lớp được bổ sung thêm $4$ học sinh mới chuyển đến nên giáo viên đã chia đều toàn bộ số học sinh trong lớp thành $4$ tổ. Hỏi hiện tại lớp học có tất cả bao nhiêu học sinh? Biết rằng số học sinh mỗi tổ theo cách chia mới ít hơn cách chia ban đầu là $2$ học sinh.
\loigiai{
Gọi số học sinh của lớp lúc đầu là $x$ (học sinh, điều kiện $x \in \mathbb{N}^*$, $x$ chia hết cho 3).
Số học sinh mỗi tổ dự định ban đầu là: $\dfrac{x}{3}$ (học sinh).
Sau khi nhận thêm $4$ học sinh, số học sinh của lớp là: $x + 4$ (học sinh).
Số học sinh mỗi tổ theo cách chia mới là: $\dfrac{x + 4}{4}$ (học sinh).
Theo đề bài, mỗi tổ hiện nay ít hơn lúc đầu $2$ học sinh nên ta có phương trình:
$\dfrac{x}{3} - \dfrac{x + 4}{4} = 2 \iff 4x - 3(x + 4) = 24 \iff x - 12 = 24 \iff x = 36$ (thỏa mãn điều kiện).
Số học sinh hiện tại của lớp là:
$36 + 4 = 40$ (học sinh).
Vậy hiện tại lớp có tất cả $40$ học sinh.
}
\end{bt}

\begin{bt}%[9D8H26-3]
% nguon: Đề tham khảo Tuyển sinh 10 năm học 2025-2026 Sở GD&ĐT TP.HCM – Quận Tân Bình - Đề 2 (Bài 6)
Thống kê doanh số bán máy tính để bàn trong $6$ tháng đầu năm của một cửa hàng thiết bị tin học ghi nhận: có đúng $4$ tháng cửa hàng bán được không quá $30$ máy tính để bàn (gồm các tháng 1, 2, 4, 6), còn $2$ tháng còn lại bán được trên $30$ máy. Đối với máy tính xách tay, có đúng $3$ tháng cửa hàng bán được từ $50$ máy trở lên (gồm các tháng 3, 5, 6). Chọn ngẫu nhiên một tháng trong $6$ tháng đầu năm.
\begin{enumerate}[a)]
    \item Tính xác suất của biến cố $A$: ``Tháng được chọn có số lượng máy tính để bàn bán được không quá $30$ chiếc''.
    \item Tính xác suất của biến cố $B$: ``Tháng được chọn có số lượng máy tính xách tay bán được ít nhất $50$ chiếc''.
\end{enumerate}
\loigiai{
Số tháng trong nửa đầu năm là $6$ tháng, do đó số kết quả có thể xảy ra khi chọn ngẫu nhiên một tháng là $n(\Omega) = 6$.
\begin{enumerate}[a)]
    \item Có $4$ tháng cửa hàng bán được không quá $30$ máy tính để bàn, do đó số kết quả thuận lợi cho biến cố $A$ là $n(A) = 4$.
    Xác suất của biến cố $A$ là:
    $P(A) = \dfrac{n(A)}{n(\Omega)} = \dfrac{4}{6} = \dfrac{2}{3}$.
    \item Có $3$ tháng cửa hàng bán được ít nhất $50$ máy tính xách tay, do đó số kết quả thuận lợi cho biến cố $B$ là $n(B) = 3$.
    Xác suất của biến cố $B$ là:
    $P(B) = \dfrac{n(B)}{n(\Omega)} = \dfrac{3}{6} = \dfrac{1}{2} = 0{,}5$.
\end{enumerate}
}
\end{bt}

\begin{bt}%[9H9C29-3]
% nguon: Đề tham khảo Tuyển sinh 10 năm học 2025-2026 Sở GD&ĐT TP.HCM – Quận Tân Bình - Đề 2 (Bài 7)
Cho tam giác nhọn $ABC$ có $\widehat{BAC} = 45^\circ$ nội tiếp đường tròn $(O; R)$. Hai đường cao $BD$ và $CE$ của tam giác cắt nhau tại trực tâm $H$.
\begin{enumerate}[a)]
    \item Chứng minh tứ giác $BEDC$ là tứ giác nội tiếp và xác định vị trí tâm $I$ của đường tròn ngoại tiếp tứ giác đó.
    \item Chứng minh $AD \cdot AC = AE \cdot AB$ và tính tỉ số $\dfrac{DE}{BC}$.
    \item Tính bán kính đường tròn ngoại tiếp tam giác $AED$ theo $R$.
\end{enumerate}
\loigiai{
\begin{enumerate}[a)]
    \item Ta có $BD \perp AC$ tại $D \implies \widehat{BDC} = 90^\circ$ và $CE \perp AB$ tại $E \implies \widehat{BEC} = 90^\circ$.
    Hai đỉnh $D$ và $E$ cùng nhìn đoạn thẳng $BC$ dưới một góc vuông, do đó tứ giác $BEDC$ nội tiếp đường tròn đường kính $BC$.
    Tâm $I$ của đường tròn ngoại tiếp tứ giác $BEDC$ là trung điểm của cạnh $BC$.
    \item Xét hai tam giác vuông $ADB$ (tại $D$) và $AEC$ (tại $E$), ta có góc $\widehat{BAC}$ chung, suy ra $\triangle ADB \backsim \triangle AEC$ (g.g).
    Do đó $\dfrac{AD}{AE} = \dfrac{AB}{AC} \implies AD \cdot AC = AE \cdot AB$.
    Từ tỉ lệ thức trên suy ra $\dfrac{AD}{AB} = \dfrac{AE}{AC}$.
    Xét $\triangle ADE$ và $\triangle ABC$, ta có góc $\widehat{BAC}$ chung và $\dfrac{AD}{AB} = \dfrac{AE}{AC}$, suy ra $\triangle ADE \backsim \triangle ABC$ (c.g.c).
    Tỉ số đồng dạng giữa hai tam giác là:
    $\dfrac{DE}{BC} = \dfrac{AD}{AB} = \cos \widehat{BAC} = \cos 45^\circ = \dfrac{\sqrt{2}}{2}$.
    \item Gọi $R'$ là bán kính đường tròn ngoại tiếp tam giác $ADE$.
    Theo định lí sin trong tam giác $ADE$:
    $2R' = \dfrac{DE}{\sin \widehat{DAE}} = \dfrac{DE}{\sin 45^\circ} = \dfrac{DE}{\dfrac{\sqrt{2}}{2}} = DE\sqrt{2}$.
    Theo định lí sin trong tam giác $ABC$:
    $2R = \dfrac{BC}{\sin \widehat{BAC}} = \dfrac{BC}{\sin 45^\circ} = BC\sqrt{2} \implies BC = R\sqrt{2}$.
    Mà $\dfrac{DE}{BC} = \dfrac{\sqrt{2}}{2} \implies DE = BC \cdot \dfrac{\sqrt{2}}{2} = R\sqrt{2} \cdot \dfrac{\sqrt{2}}{2} = R$.
    Do đó $2R' = DE\sqrt{2} = R\sqrt{2} \implies R' = \dfrac{R\sqrt{2}}{2}$.
    Vậy bán kính đường tròn ngoại tiếp tam giác $AED$ là $R' = \dfrac{R\sqrt{2}}{2}$.
\end{enumerate}
}
\end{bt}
